% =====================================================================
%  paper_summary_v1.tex  --  structure of the paper, one summary per section
%  Started 2026-11-17 from paper/plan_v0.tex and the fact sheets in
%  paper/factsheets/ (every number below is taken from them; the fact
%  sheets give the run directory / csv each number comes from).
%
%  Colour code used in the text:
%    red    [TODO: ...]   something to run, check or decide
%    blue   [note]        comment for the author, not paper text
%    orange numbers       PROVISIONAL: study II numbers, to be replaced by
%                         the WSD redo (one batch size, no early stopping)
%  Figure boxes are placeholders; the path printed in each box is relative
%  to /storage3/DSIP/rriva/research/ and points to the existing plot.
%  Compile with pdflatex + bibtex (refs.bib in the same folder).
% =====================================================================
\documentclass[11pt,a4paper]{article}

\usepackage[T1]{fontenc}
\usepackage{lmodern}
\usepackage[a4paper,top=2.2cm,bottom=2.2cm,left=2.2cm,right=2.2cm]{geometry}
\usepackage{amsmath,amssymb}
\usepackage{graphicx}
\usepackage{booktabs}
\usepackage{multirow}
\usepackage{array}
\usepackage{xcolor}
\usepackage{url}
\usepackage[numbers,sort&compress]{natbib}
\usepackage{caption}
\usepackage[colorlinks=true,linkcolor=blue!50!black,citecolor=green!40!black,urlcolor=blue!50!black]{hyperref}

\captionsetup{font=small,labelfont=bf}

% ---------------------------------------------------------------- macros
\newcommand{\todo}[1]{\textcolor{red}{\textbf{[TODO:} #1\textbf{]}}}
\newcommand{\note}[1]{\textcolor{blue!70!black}{[\emph{#1}]}}
\newcommand{\prov}[1]{\textcolor{orange!85!black}{#1}}
\newcommand{\LCE}{\mathcal{L}_{\mathrm{CE}}}
\newcommand{\LJC}{\mathcal{L}_{\mathrm{JC}}}
\newcommand{\Ltrk}{\mathcal{L}_{\mathrm{trk}}}
\newcommand{\Ljet}{\mathcal{L}_{\mathrm{jet}}}
\newcommand{\from}[2]{#1$\leftarrow$#2}
\newcommand{\ttbar}{\ensuremath{t\bar{t}}}
\newcommand{\pT}{\ensuremath{p_{\mathrm{T}}}}

% placeholder for a figure that already exists as a plot
\newcommand{\figph}[4]{%
  \begin{figure}[htbp]
    \centering
    \fbox{\parbox[c][#1][c]{0.93\linewidth}{\centering\small\textit{Figure placeholder}\\[4pt]
      \footnotesize existing plot: \url{#2}}}
    \caption{#3}
    \label{#4}
  \end{figure}}
% placeholder for a figure that still has to be made
\newcommand{\figtodo}[4]{%
  \begin{figure}[htbp]
    \centering
    \fbox{\parbox[c][#1][c]{0.93\linewidth}{\centering\small\textit{Figure placeholder (to be made)}\\[4pt]
      \footnotesize #2}}
    \caption{#3}
    \label{#4}
  \end{figure}}

% ------------------------------------------------------------- front matter
\title{Contrastive self-supervision for jet flavour tagging:\\
augmentation strength, intersample attention and training strategies\\[4pt]
\large Summary draft v1 (structure of the paper, one summary per section)}
\author{Riccardo Riva, Marco Cristoforetti, Roberto Iuppa, Francesco Maria Follega
\todo{confirm author list, order and affiliations (FBK, Univ.\ Trento, TIFPA?)}}
\date{17 November 2026}

\begin{document}
\maketitle

% =====================================================================
\begin{abstract}
Deep-learning flavour taggers compress the features of a jet and of its tracks into a
latent embedding from which the jet flavour is predicted. We study whether contrastive
self-supervised objectives acting directly on this embedding improve flavour tagging,
using the public ATLAS Open Data JetSet sample (simulated \ttbar{} events at
$\sqrt{s}=13.6$\,TeV). Jets are treated as tabular data and encoded by a SAINT-like
transformer with per-feature tokens, attention across features and intersample attention
across the jets and tracks of a batch, followed by a GN2-like track aggregator. An InfoNCE
loss between an augmented and a clean view of each jet is used in two ways:
(i) multi-task training, where it is added to the classification loss, and
(ii) contrastive pretraining on a disjoint sample of 40M jets, followed by training from
scratch, fine-tuning or training on frozen encoders. In multi-task training the augmentation
strength is the decisive parameter. Weak augmentations degrade the classifier. Strong ones
delay overfitting, spread the jet embedding uniformly over the hypersphere and give the
lowest validation loss. They also give a c-jet rejection about 2\% higher than a GN2-like
tagger trained on the same jets, at equal light- and $\tau$-jet rejection. The same objective
has no effect on the GN2-like tagger. In the pretraining study, only pretraining the full
network including the aggregator, followed by fine-tuning, outperforms training from scratch:
\prov{+11\%} c-jet rejection at equal light- and $\tau$-jet rejection. Pretraining only the
per-feature embeddings and freezing the encoders do not help, at any fraction of labelled
jets between 1\% and 100\%. We also show that contrastive pretraining with intersample
attention needs query--key normalisation to avoid unbounded attention-logit growth. Finally,
we discuss embedding-quality metrics (uniformity, effective rank) as diagnostics rather than
targets.
\end{abstract}

\noindent\note{Abstract = the whole paper in 180--250 words. Rewrite it last, once the
numbers are final. The study-II numbers (orange) will change with the WSD redo.}

\medskip
\noindent\todo{DECISION before writing the results: the two studies currently use different
models and datasets (study I: Model A, 2.2M parameters, 52M training jets; study II: Model B,
7.5M parameters, 18.6M training jets). Either (a) present them as two explicitly separate
set-ups, or (b) add a CE+JC $\lambda$ ladder of Model B and a GN2 CE / CE+JC pair on the
same 18.6M-jet split to the WSD redo, so the whole paper uses one model and one dataset,
with the Model-A scan as supporting material. See the review in the chat of 2026-11-17.}

\tableofcontents
\newpage

% =====================================================================
\section{Introduction}
\label{sec:intro}

\note{Purpose: why flavour tagging matters, what is missing in current taggers (an
explicitly shaped embedding), what we do, and the list of contributions. Write this section
last; the contributions list is the paper's spine and should be written first.}

\paragraph{Flavour tagging.}
Identifying jets that originate from $b$- and $c$-quarks (and from hadronic $\tau$ decays)
is central to the LHC physics programme: Higgs couplings to heavy quarks, top physics and
many searches. It relies on the long lifetime of $b$- and $c$-hadrons, which shows up as
tracks with large impact parameters and as displaced vertices inside the jet.

\paragraph{Transformer taggers and their embedding.}
Flavour taggers have moved from likelihood and BDT combinations of hand-crafted
variables~\citep{ATLAS:2019} to graph and transformer networks that read the tracks
directly: GN1~\citep{GN1} and GN2~\citep{GN2} in ATLAS, and DeepJet~\citep{Bols:2020},
ParticleNet~\citep{Qu_2020}, ParT~\citep{Qu:2022} and UParT~\citep{CMSDP2025081} in CMS.
All of them map the jet to a latent representation from which the flavour is predicted.
That representation drives the performance, and possibly the behaviour outside the
training sample, yet it is shaped only indirectly, through the classification loss and
auxiliary tasks.

\paragraph{Jets as tabular data.}
The inputs of a jet (two jet-level kinematic features and, for each track, 19 heterogeneous
continuous and categorical features) form a small table. Deep tabular models such as
SAINT~\citep{somepalli2021saint} embed every feature individually and combine attention
across features with \emph{intersample attention} across the rows of a batch. GN-like
taggers instead embed the whole feature vector of a track at once. We adapt the tabular
approach to jets and tracks.

\paragraph{Contrastive self-supervision.}
Contrastive learning shapes an embedding directly. An InfoNCE loss~\citep{oord2018representation,Chen2020SimCLR}
pulls together two views of the same object and pushes apart other objects. It optimises
the alignment and uniformity of the embedding~\citep{Wang2020AlignmentUniformity}, and the
strength of the augmentation controls how much information the views share~\citep{tian2020makes}.
In jet physics, contrastive learning has mostly used physics-motivated augmentations,
for symmetries and safety~\citep{Dillon2022JetCLR}, re-simulation~\citep{Harris2025Resim}
or data/simulation differences~\citep{Sheldon2025MACK}. Here we use generic,
feature-agnostic tabular augmentations (CutMix and Mixup with strength $\lambda$).

\paragraph{This work.}
We ask two questions:
(Q1) does adding a contrastive term on the jet embedding to the classification loss
(\emph{multi-task training}) improve a tagger, and how does the answer depend on the
augmentation strength and on the architecture?
(Q2) does \emph{contrastive pretraining} on a large unlabelled sample help, which
pretraining objective works, and is it better used by fine-tuning or with frozen encoders?
For both questions we report flavour-tagging performance and embedding-quality metrics.

\paragraph{Contributions} \note{draft, to be sharpened once the results are final}
\begin{itemize}
  \item A SAINT-like jet tagger with per-feature tokens, column attention and intersample
        attention across the jets and tracks of a batch, combined with a GN2-like track
        aggregator (Section~\ref{sec:model}).
  \item A scan of the augmentation strength $\lambda$ in multi-task training (cross-entropy
        plus jet-embedding contrastive loss), repeated for a GN2-like tagger trained on the
        same jets: the contrastive term helps our model only for strong augmentations and
        is neutral for GN2 (Section~\ref{sec:res-multitask}).
  \item A controlled comparison of contrastive pretraining recipes (P0, P1) and downstream
        strategies (from scratch, fine-tuning, frozen encoders) on a 40M-jet pretraining
        split, including a scan of the labelled-data fraction (Section~\ref{sec:res-pretrain}).
  \item Training-stability findings for intersample attention under contrastive objectives:
        attention-logit growth cured by QK-normalisation, and collapse of multi-task training
        from scratch without tight gradient clipping (Section~\ref{sec:stability}).
\end{itemize}

\paragraph{Outline.} Sec.~\ref{sec:related} related work; Sec.~\ref{sec:data} dataset;
Sec.~\ref{sec:model} model; Sec.~\ref{sec:methods} objectives, training strategies and
metrics; Secs.~\ref{sec:res-multitask}--\ref{sec:stability} results;
Sec.~\ref{sec:discussion} discussion and limitations; Sec.~\ref{sec:conclusions} conclusions.

% =====================================================================
\section{Related work}
\label{sec:related}

\note{Purpose: place the work in four literatures in about 1.5 pages. Each paragraph ends
with how we differ. The "Background" part of plan\_v0 (foundation models, contrastive
loss, uniformity, effective rank) moves to Section~\ref{sec:methods}, where the
quantities are defined and used.}

\subsection{Transformer-based flavour tagging}
GN1/GN2 in ATLAS~\citep{GN1,GN2} (implemented in Salt~\citep{Barr2025Salt}) and ParT/UParT
in CMS~\citep{Qu:2022,CMSDP2025081}: per-track embeddings, self-attention over the tracks
of one jet, attention pooling, auxiliary track-origin and vertexing tasks. The GN2 paper
also released the JetSet open-data sample~\citep{ATLASJetSet2025} we use. \emph{Difference:}
no existing tagger shapes its embedding with a self-supervised loss, and none lets a jet
attend to other jets.

\subsection{Deep learning for tabular data and intersample attention}
Gradient-boosted trees remain strong on tabular data~\citep{Grinsztajn2022TreesTabular}.
Attention models close the gap with per-feature tokens: TabTransformer~\citep{Huang2020TabTransformer},
FT-Transformer~\citep{Gorishniy2021FTTransformer}, and SAINT~\citep{somepalli2021saint},
which adds intersample (row) attention and contrastive pretraining. Non-parametric
transformers~\citep{Kossen2021NPT} and the Set Transformer's induced attention~\citep{lee2019settransformerframeworkattentionbased}
are related ways of attending across data points. Tabular foundation models (TabPFN~\citep{Hollmann2025TabPFNNature},
TabICL~\citep{Qu2025TabICL}) are the long-term motivation. \emph{Difference:} our
"table" is hierarchical (jets and their variable-length track sets), and we apply row
attention to both levels.

\subsection{Contrastive self-supervised learning}
InfoNCE / CPC~\citep{oord2018representation}, SimCLR~\citep{Chen2020SimCLR},
MoCo~\citep{He2020MoCo}, SupCon~\citep{Khosla2020SupCon}. Alignment and
uniformity~\citep{Wang2020AlignmentUniformity}; views and InfoMin~\citep{tian2020makes};
tabular augmentations: SCARF~\citep{Bahri2022SCARF}, VIME~\citep{Yoon2020VIME}, the
CutMix/Mixup corruption of SAINT~\citep{Yun2019CutMix,Zhang2018mixup};
dimensional collapse and effective rank~\citep{Jing2022DimCollapse,Roy2007EffectiveRank,Garrido2023RankMe};
contrastive terms added to supervised training~\citep{Gunel2021SCLFinetune}; fine-tuning
versus linear probing~\citep{Kumar2022FTDistorts}. \emph{Difference:} we measure how
the augmentation strength interacts with a supervised loss on a physics task.

\subsection{Self-supervision and foundation models in jet physics}
JetCLR~\citep{Dillon2022JetCLR}; re-simulation-based SSL (RS3L)~\citep{Harris2025Resim};
masked particle modelling~\citep{golling2024masked}; OmniJet-$\alpha$~\citep{Birk2024OmniJet};
OmniLearn~\citep{Mikuni2025OmniLearn}; fine-tuning foundation models for
analyses~\citep{vigl2024finetuning}; contrastive approaches to mismodelling~\citep{Sheldon2025MACK}.
\emph{Difference:} generic tabular augmentations, a flavour-tagging target, intersample
attention, and an explicit comparison of multi-task training with pretraining.
\todo{check whether any SSL work targets flavour tagging specifically; the literature
fact sheet lists 2025--26 entries (HEP-JEPA, LHCb SSL, scaling laws) that are worth one
sentence each.}

% =====================================================================
\section{Dataset}
\label{sec:data}

\note{Purpose: everything needed to reproduce the inputs. Short, factual, with one
table of the splits and one of the input features.}

\subsection{JetSet sample and selection}
ATLAS Open Data JetSet~\citep{ATLASJetSet2025}: about 48.7M simulated \ttbar{} events at
$\sqrt{s}=13.6$\,TeV, 199.8M jets with up to 40 associated tracks. Jets are required to
have $20<\pT<250$\,GeV and $|\eta|<2.5$, which leaves about 192M jets. Flavour labels
($b$, $c$, light, $\tau$) come from the truth hadron/lepton cone labelling.
\todo{check jet collection and track association against the Open Data record; check the
13.6\,TeV value in the record (some old plot labels say 13\,TeV).}

\subsection{Preprocessing and splits}
Preprocessing with UPP: \pT--$\eta$ resampling (``countup'') to the $c$-jet spectrum on a
$50\times40$ grid, fixed flavour composition $b:c:\text{light}:\tau = 25.8:25.8:38.7:9.7\%$,
and $z$-score normalisation with the training-split statistics. Splits are defined by the
event number, so no event is shared between splits (Table~\ref{tab:splits}). Study I uses
dataset~A. Study II uses dataset~B, whose ``pretrain'' split serves for contrastive
pretraining with the labels unused.

\begin{table}[htbp]
\centering\small
\caption{Dataset splits. Both datasets come from the same JetSet events with the same
preprocessing; 43\% of the distinct jets of every dataset-B split are also in the dataset-A
training split, so models trained on A are never evaluated on B.}
\label{tab:splits}
\begin{tabular}{llllll}
\toprule
Dataset & Split & event-number rule & Jets & Used for \\
\midrule
A (\texttt{251112\_112022}) & train & $\bmod 62 \in [0,51]$  & 52.0M  & Study I training \\
                            & val   & $\bmod 62 \in [52,56]$ & 5.0M   & early stopping, model selection \\
                            & test  & $\bmod 62 \in [57,61]$ & 5.0M   & performance \\
\midrule
B (\texttt{251112\_152051}) & pretrain & $\bmod 100 \in [0,63]$  & 39.68M & contrastive pretraining (no labels) \\
                            & train    & $\bmod 100 \in [64,93]$ & 18.60M & Study II supervised training \\
                            & val      & $\bmod 100 \in [94,96]$ & 1.86M  & early stopping, fraction search \\
                            & test     & $\bmod 100 \in [97,99]$ & 1.86M  & performance \\
\bottomrule
\end{tabular}
\end{table}

\note{Say explicitly that the ``unlabelled'' pretraining split still has the flavour
composition fixed by truth-based resampling, so it is not a realistic unlabelled sample
(limitation, Section~\ref{sec:discussion}).}

\subsection{Input features}
Two jet features (\texttt{pt\_btagJes}, \texttt{eta\_btagJes}) and 19 track features: 10
continuous (\texttt{d0}, \texttt{z0SinTheta}, \texttt{dphi}, \texttt{deta}, \texttt{qOverP},
the two lifetime-signed impact-parameter significances and three uncertainties) and 9
categorical hit counts (pixel, SCT, innermost-layer, shared and split hits). These are the
same inputs as the GN2-like reference. On average a jet has 7.15 valid tracks (82\% of the
40-slot grid is padding); $\tau$-jets have 3.7.
\todo{full feature table in an appendix, with the embedding cardinality of each
categorical feature.}

% =====================================================================
\section{Model}
\label{sec:model}

\note{Purpose: describe the architecture once, then the two variants used in the two
studies, then the GN2 reference. The architecture figure does not exist yet
(\texttt{images/fm4tag.png} in plan\_v0 was never drawn); \texttt{diary/figures/gn2.png}
is the ATLAS figure with handwritten notes and must be redrawn or cited.}

\subsection{Per-feature tokenisation}
Every input feature becomes its own 16-dimensional token. Each continuous feature has its
own small MLP ($1\to32\to16$), and the categorical hit counts share one embedding table
with per-feature offsets. A jet is a $2\times16$ token grid and a track a $19\times16$ grid.

\subsection{Column and intersample (row) attention}
The encoder is a stack of depth 6. Each step applies pre-norm \emph{column} attention
across the feature tokens of one object, then pre-norm \emph{row} attention across
objects: jets attend to other jets, and tracks attend to the valid tracks of
\emph{all} jets. Feed-forward layers use GEGLU~\citep{Shazeer2020GLU}, and the residual
is a true pre-norm, $x+f(\mathrm{LN}(x))$~\citep{Xiong2020PreLN}. A per-object projector
maps the tokens to per-track embeddings (128-d) and a jet-feature embedding (32-d in
Model A, 128-d in Model B).
Two variants of the row attention are used (Table~\ref{tab:variants}):
\begin{itemize}
  \item \textbf{Model A} (study I): induced set attention~\citep{lee2019settransformerframeworkattentionbased}
        with 16 inducing points per feature token (``per-token'' mode) over the whole
        batch on each GPU.
  \item \textbf{Model B} (study II): SAINT-style ``concat'' mode, where the feature tokens
        of an object are concatenated into one vector, with full softmax attention inside
        chunks of 512 consecutive jets (variable-length FlashAttention~\citep{Dao2024FlashAttention2})
        and RMS normalisation of queries and keys~\citep{Henry2020QKNorm,Zhang2019RMSNorm}
        (Section~\ref{sec:stability}).
\end{itemize}

\subsection{Track aggregator and classification head}
A 6-layer transformer runs over the track embeddings of each jet separately (128-d,
$4\times64$ heads), followed by a mean pool. The pooled vector is concatenated with the
jet-feature embedding to form the \emph{global jet embedding} $z_{\mathrm{jet}}$ (160-d in
Model A, 256-d in Model B). An MLP head ($z_{\mathrm{jet}}\to256\to4$) predicts the
flavour. There are no auxiliary tasks.

\begin{table}[htbp]
\centering\small
\caption{The two model variants and the GN2-like reference.}
\label{tab:variants}
\begin{tabular}{p{4.2cm}p{3.4cm}p{3.6cm}p{3.6cm}}
\toprule
 & Model A (study I) & Model B (study II) & GN2-like (Salt) \\
\midrule
Feature embedding      & per feature, 16-d & per feature, 16-d & per track (init MLP 256) \\
Row (intersample) attention & ISAB, 16 inducing points, per token, whole batch & full attention in 512-jet chunks, concat, QK-norm & none \\
Track transformer      & aggregator, 6 layers & aggregator, 6 layers & 4 layers, 256-d \\
Pooling                & mean & mean & attention pooling \\
Jet embedding dim.     & 160 & 256 & 128 \\
Trainable parameters   & 2.22M (89\% in aggregator) & 7.47M (70\% in track row attention) & 2.24M \\
Contrastive negatives (JC) & global batch (2047 or 4095) & 511 (same chunk) & global batch (4095) \\
Dataset                & A & B & A (study I), B (reference) \\
\bottomrule
\end{tabular}
\end{table}

\subsection{GN2-like reference}
The reference is a GN2 architecture trained in Salt~\citep{Barr2025Salt} on the same inputs
and training jets, with the jet-classification task only (no track-origin or vertexing
heads). It has 2.2M parameters, about the same size as Model A. The production GN2v01
scores stored in the open-data files are shown only as an external yardstick: that tagger
was trained on a different, larger sample.
\note{Call it ``GN2-like'' everywhere: without the auxiliary tasks it is not the ATLAS GN2.}

\figtodo{5cm}{To draw: per-feature tokens $\to$ [column attention $\to$ row attention]
$\times6$ $\to$ projectors $\to$ track aggregator + mean pool $\to$ $z_{\mathrm{jet}}$
$\to$ head. Mark where $\LCE$, $\LJC$, $\Ltrk$ and $\Ljet$ act. Blueprints:
\texttt{fm4tag/README.md}, \texttt{fm4tag-chunk/docs/encoder\_dataflow.md}.}
{Architecture of the tagger, and where each loss acts.}{fig:arch}

% =====================================================================
\section{Learning objectives, training strategies and metrics}
\label{sec:methods}

\note{Purpose: define every loss, every training recipe and every metric used in
Sections~\ref{sec:res-multitask}--\ref{sec:stability}. The P0/P1 table
(Table~\ref{tab:kinds}) is the centre of this section.}

\subsection{Supervised loss}
Unweighted four-class cross-entropy $\LCE$ on the clean (un-augmented) forward pass.
\note{Salt's GN2 uses class-weighted CE by default. Every CE value in this paper must be
unweighted, or labelled as weighted.}

\subsection{Augmented views}
Two views of each batch: a clean view and a view corrupted with strength
$\lambda\in[0,1]$. \emph{CutMix} replaces each normalised feature value, with probability
$\lambda$, by the same feature of a randomly chosen other row (another jet, or another
track of the batch). \emph{Mixup} then sets $x'=(1-\lambda)\,x+\lambda\,x_{\pi}$ on the
feature embeddings. One $\lambda$ drives both. For $\lambda\ge0.5$ the Mixup view is closer
to its partner than to itself.

\subsection{Contrastive losses}
All contrastive terms use the same symmetric cross-view InfoNCE on L2-normalised embeddings
with temperature $\tau=0.1$. The positive is the other view of the same object; the
negatives are the other objects of the other view. We use two such terms:
\begin{itemize}
  \item the \textbf{jet-embedding contrastive loss} $\LJC$, on the global jet embedding
        $z_{\mathrm{jet}}$ (the input of the head);
  \item the \textbf{per-object losses} $\Ltrk$ and $\Ljet$, on the per-track and the
        jet-feature embeddings at the encoder output (pretraining only).
\end{itemize}
\note{A denoising/reconstruction loss is implemented but used in no reported run. Remove it
from the paper, and from the poster wording ``and a denoising reconstruction loss''.}

\subsection{Multi-task training (study I)}
$\mathcal{L}=\LCE+w\,\LJC$ with $w=0.3$, trained from scratch. The CE-only baseline is
$w=0$, where $\lambda$ has no effect.

\subsection{Contrastive pretraining and downstream training (study II)}
Pretraining runs on the pretrain split with no labels in the loss, with $\lambda=0.4$ and
$\tau=0.1$. These were the best values of a pilot grid judged by a frozen-encoder probe
(Appendix~\ref{app:pilots}). The recipes are listed in Table~\ref{tab:pretrain}.
Downstream, the tagger is trained on the labelled train split in one of three ways:
\begin{itemize}
  \item \textbf{SC}, from scratch;
  \item \textbf{FT}, fine-tuning: start from the pretrained weights, keep them fixed for 2
        epochs, then train them at 1/10 of the learning rate;
  \item \textbf{FR}, frozen encoders: encoders and projectors frozen; aggregator and head
        trained.
\end{itemize}
Each method is run with or without the downstream term $0.3\,\LJC$.
Tables~\ref{tab:kinds} and~\ref{tab:downstream} list every combination that was run.

\begin{table}[htbp]
\centering\small
\caption{Contrastive pretraining recipes (study II, Model B, 2 seeds each).
P2 was run only as a pilot and is reported in Appendix~\ref{app:p2}.}
\label{tab:pretrain}
\begin{tabular}{p{3.6cm}p{3.7cm}p{3.9cm}p{3.7cm}}
\toprule
 & \textbf{P0} & \textbf{P1} & P2 (pilot only) \\
\midrule
Loss & $\Ltrk+\Ljet$ & $\Ltrk+\Ljet+\LJC$ & $\LJC$ \\
Acts on & track and jet-feature embeddings & + global jet embedding & global jet embedding only \\
Trained modules & encoders, projectors & encoders, projectors, aggregator & encoders, projectors, aggregator \\
Aggregator after pretraining & untrained (random) & pretrained & pretrained \\
\midrule
Data & \multicolumn{3}{p{11.9cm}}{pretrain split, 39.68M jets, labels not used} \\
Views and loss & \multicolumn{3}{p{11.9cm}}{CutMix + Mixup ($\lambda=0.4$) vs.\ clean; $\tau=0.1$; negatives within the 512-jet chunk} \\
Optimisation & \multicolumn{3}{p{11.9cm}}{AdamW, peak lr $10^{-3}$, 2\% warm-up + cosine, 30 epochs, batch 12\,288, clip 1.0, QK-norm} \\
\midrule
Final val.\ loss (two seeds) & 5.44 / 5.51 & 5.67 / 5.73 & 0.222 / 0.223 \\
Eff.\ rank: tracks / jet feat. / $z_{\mathrm{jet}}$ & 121 / 62 / -- & 126 / 62 / 130 & 59 / 1 / 127 \\
Wall time per seed & 43\,h (A40) & 34\,h (L40S) & 43\,h (A40) \\
Used downstream & FT, FR & FT, FR, with and without JC & no \\
\bottomrule
\end{tabular}
\end{table}

\begin{table}[htbp]
\centering\small
\caption{Kinds of training done with P0 and P1 (study II). \checkmark: run with 4 seeds
(pretrained arms: 2 pretraining $\times$ 2 downstream seeds). The downstream loss is
$\LCE$ or $\LCE+0.3\,\LJC$.}
\label{tab:kinds}
\begin{tabular}{l|cc|cc|cc}
\toprule
 & \multicolumn{2}{c|}{From scratch (SC)} & \multicolumn{2}{c|}{Fine-tuning (FT)} & \multicolumn{2}{c}{Frozen encoders (FR)} \\
Pretraining & CE & CE+JC & CE & CE+JC & CE & CE+JC \\
\midrule
none & \checkmark & \checkmark & -- & -- & -- & -- \\
P0   & -- & -- & \checkmark & not run$^{a}$ & \checkmark & not run$^{a}$ \\
P1   & -- & -- & \checkmark & \checkmark & \checkmark & \checkmark \\
P2   & -- & -- & not run & not run & probe only$^{b}$ & not run \\
\bottomrule
\multicolumn{7}{l}{\footnotesize $^{a}$ by design: P0 leaves the aggregator untrained.
$^{b}$ 15-epoch frozen probe in the pilot grid.}
\end{tabular}
\end{table}

\begin{table}[htbp]
\centering\small
\caption{Downstream training recipes (study II). Common settings: train split (18.6M jets),
AdamW, 2\% warm-up + cosine over 40 epochs, early stopping (patience 12) on the validation
CE, JC with $w=0.3$, $\lambda=0.4$, $\tau=0.1$. FT holds the pretrained parts at lr\,=\,0
for 2 epochs. \prov{To be replaced by the WSD redo: one batch size (12\,800) for every
method, warm-up--stable--decay schedule, no early stopping.}}
\label{tab:downstream}
\begin{tabular}{p{2.3cm}p{3.2cm}p{3.0cm}p{1.5cm}p{1.3cm}p{2.7cm}}
\toprule
Method & Initialisation & Trainable & Loss & Batch & lr (pretrained) \\
\midrule
SC              & random & all & CE & 28\,672 & $1.1\cdot10^{-3}$ \\
SC+JC           & random & all (grad.\ clip 1.0) & CE+JC & 14\,336 & $8\cdot10^{-4}$ \\
\from{FT}{P0}   & P0 encoders; aggregator random & all & CE & 28\,672 & $1.1\cdot10^{-3}$ ($1.1\cdot10^{-4}$) \\
\from{FR}{P0}   & P0 encoders; aggregator random & aggregator, head & CE & 94\,208 & $2\cdot10^{-3}$ \\
\from{FT}{P1}   & P1 encoders + aggregator & all & CE & 28\,672 & $1.1\cdot10^{-3}$ ($1.1\cdot10^{-4}$) \\
\from{FR}{P1}   & P1 encoders + aggregator & aggregator, head & CE & 94\,208 & $2\cdot10^{-3}$ \\
\from{FT}{P1}+JC & P1 encoders + aggregator & all & CE+JC & 14\,336 & $8\cdot10^{-4}$ ($8\cdot10^{-5}$) \\
\from{FR}{P1}+JC & P1 encoders + aggregator & aggregator, head & CE+JC & 53\,760 & $1.5\cdot10^{-3}$ \\
\bottomrule
\end{tabular}
\end{table}

\subsection{Embedding-quality metrics}
On $\le 8192$ validation jets per epoch:
\begin{itemize}
  \item \emph{uniformity}~\citep{Wang2020AlignmentUniformity},
        $\log\langle e^{-2\|\hat z_i-\hat z_j\|^2}\rangle$ on unit-normalised embeddings;
        lower is more uniform, about $-4$ in the high-dimensional limit;
  \item \emph{effective rank}~\citep{Roy2007EffectiveRank}, the exponential of the entropy
        of the normalised singular values of the centred embedding.
\end{itemize}
\note{Alignment, listed in plan\_v0, is not implemented. Either add it (cheap: mean
$\|\hat z-\hat z'\|^2$ between the two views) or drop it from the text.}

\subsection{Flavour-tagging performance}
The $b$-tagging discriminant is
$D_b=\ln\!\big[p_b/(f_c\,p_c+f_\tau\,p_\tau+(1-f_c-f_\tau)\,p_{\mathrm{light}})\big]$.
The rejection of flavour $X$ at a $b$-efficiency working point is $1/\varepsilon_X$. We use
two settings of the fractions:
\begin{itemize}
  \item nominal: $f_c=0.2$, $f_\tau=0.01$;
  \item \emph{constrained optimal}: per model, maximise the $c$-rejection on a grid of
        $(f_c,f_\tau)$ (step 0.002), subject to light- and $\tau$-rejection at least equal
        to the reference's.
\end{itemize}
\todo{one protocol for the whole paper: optimise the fractions on the validation split, then
evaluate on test, at one working point. Study~I currently optimises on the test split at
80\%; study~II optimises on validation at 80\% and quotes test at 77\%.}

\subsection{Training protocol and statistics}
Optimiser AdamW~\citep{Loshchilov2019AdamW}, bf16 mixed precision. Each configuration is
run with 2--8 seeds. We quote mean $\pm$ standard deviation of the best validation CE.
Selection rules are fixed before the comparison:
\begin{itemize}
  \item study I: ``stuck'' seeds (validation CE never below 0.610 within 20 epochs) are
        excluded from the scan and the baselines alike and listed;
  \item study II: collapse rule (validation CE $>0.85$ after a best epoch within the first
        6), pre-registered.
\end{itemize}
Learning rates are tuned per method on one seed (42), which is excluded from the final runs.
Compute: Appendix~\ref{app:hyper}.

% =====================================================================
\section{Results I: multi-task training and augmentation strength}
\label{sec:res-multitask}

\note{Purpose: answer Q1. Claim: strong augmentation turns the contrastive term into a
useful regulariser for our model (Model A) but not for GN2. All numbers are on dataset A
(52M training jets), with the unweighted validation CE on the 5M-jet validation split.
This is the plan\_v0 2$\times$2 grid \{GN2, ours\} $\times$ \{CE, CE+CL\}.}

\subsection{Augmentation-strength scan with Model A}
We scan $\lambda$ from 0.05 to 0.8 at $w=0.3$, $\tau=0.1$, lr $3\cdot10^{-4}$; 67 trainings,
of which 6 are excluded as stuck (Table~\ref{tab:lambda-fm4tag}, Fig.~\ref{fig:lambda-fm4tag}).
\begin{itemize}
  \item Weak augmentations ($\lambda\le0.2$) make the validation CE \emph{worse} than CE
        only ($0.60345\pm0.00049$).
  \item The loss improves with $\lambda$ (not strictly: $\lambda=0.5$ is an outlier) and
        reaches a plateau at
        $\lambda\approx0.55$--$0.65$ ($\approx0.5999$), $0.0035$ below CE only.
  \item Stronger augmentation delays overfitting. The best epoch moves from about 8
        (CE only) to 60--99, and the validation CE barely rises after its minimum.
  \item Stuck seeds: 3 of 8 for CE only, 3 of 52 with JC.
\end{itemize}
Regime-pure check: in the same batch regime and the same 50-epoch schedule, $\lambda=0.3$
reaches $0.60192\pm0.00026$ against $0.60345\pm0.00049$ for CE only ($n=4$ each).

\begin{table}[htbp]
\centering\small
\caption{Model A, multi-task $\LCE+0.3\,\LJC$: best validation CE (unweighted, mean $\pm$
std over $n$ seeds) and change of the constrained-optimal $c$-rejection at 80\% $b$-efficiency
with respect to the GN2-like reference (best seed per row). Batch regimes: A = $1024\times2$
GPUs, B = $2048\times1$, C = $2048\times2$. Rows with one usable seed
($\lambda=0.64,0.7,0.75,0.8$: 0.5992--0.6001, +1.3\% to +2.5\%) are omitted.}
\label{tab:lambda-fm4tag}
\begin{tabular}{lllccc}
\toprule
$\lambda$ & regime & $n$ & best val.\ CE & best epochs & $\Delta$ c-rej vs GN2-like \\
\midrule
CE only & A & 4 & $0.60345\pm0.00049$ & 5--10   & $-2.2\%$ \\
0.05    & A & 6 & $0.60742\pm0.00047$ & 7--11   & $-5.6\%$ \\
0.1     & A & 5 & $0.60620\pm0.00035$ & 9--14   & $-3.6\%$ \\
0.2     & A, B & 3 & $0.60461\pm0.00019$ & 17   & $-3.0\%$ \\
0.25    & A & 4 & $0.60243\pm0.00023$ & 17--20  & $-1.3\%$ \\
0.3     & A, B & 5 & $0.60206\pm0.00038$ & 20--27 & $+0.4\%$ \\
0.35    & A, B & 5 & $0.60207\pm0.00059$ & 22--32 & $+0.4\%$ \\
0.4     & B, C & 3 & $0.60179\pm0.00047$ & 27--48 & $+0.0\%$ \\
0.45    & B, C & 2 & $0.60145\pm0.00041$ & 48, 72 & $+0.6\%$ \\
0.5     & A, B, C & 5 & $0.60231\pm0.00032$ & 36--46 & $-0.7\%$ \\
0.55    & B, C & 2 & $0.59987\pm0.00081$ & 62, 98 & $+2.2\%$ \\
0.6     & A, B, C & 3 & $0.59993\pm0.00078$ & 48--95 & $+2.1\%$ \\
0.65    & B, C & 2 & $0.59973\pm0.00053$ & 62, 71 & $+2.1\%$ \\
\bottomrule
\end{tabular}
\end{table}

\figph{4.5cm}{fm4tag/plots/val_loss_ce/val_loss_ce_best_vs_lambda_no_dropout_min2_categorical.png}
{Model A: mean best validation CE versus $\lambda$, with the seed spread (seeds with
$\ge2$ usable runs). Companion: per-seed curves with the best epoch marked
(\texttt{...\_categorical\_curves.png}).}{fig:lambda-fm4tag}

\todo{Controls a referee will ask for. (1) $\lambda$ is confounded with the batch regime and
the schedule: low-$\lambda$ runs stop early at the peak learning rate, high-$\lambda$ runs
reach a learning rate near zero. Across healthy JC runs, Spearman(lr at best epoch, best CE)
$=+0.91$. Needed: a regime-pure re-run (one batch size, fixed-length schedule or WSD, no early
stopping) with CE only, a dropout baseline and 4--5 values of $\lambda$. (2) The CE-only +
dropout-0.1 baseline ($0.60162\pm0.00096$, regime C) matches $\lambda=0.3$--$0.5$, so only
$\lambda\ge0.55$ beats a simple regulariser. Discuss this openly. (3) More seeds at
$\lambda\ge0.45$ ($n=1$--$2$ now).}

\subsection{The same study for the GN2-like tagger}
Same recipe in Salt ($w=0.3$, $\tau=0.1$, unweighted CE, lr $3\cdot10^{-4}$,
$2\times2048$; 41 seeds, no divergence; Table~\ref{tab:lambda-gn2}):
\begin{itemize}
  \item the contrastive term is \emph{neutral} for $\lambda\le0.4$ and slightly worse and
        noisier for $\lambda\ge0.5$;
  \item it still delays the best epoch (16 $\to$ 83);
  \item the test $c$-rejection at 77\% varies by less than 1\% across $\lambda$; the light
        and $\tau$ rejections scatter by up to about 10\%, with no trend in $\lambda$;
  \item the InfoNCE lower bound on the mutual information between the views,
        $I_{\mathrm{NCE}}=\log K-\mathcal{L}_{\mathrm{NCE}}$, falls monotonically with
        $\lambda$, which makes $\lambda$ a clean knob on view information (InfoMin~\citep{tian2020makes}).
\end{itemize}
Hyper-parameter searches (Optuna~\citep{Akiba2019Optuna}, seven studies) over $w$, $\tau$,
$\lambda$ and the learning rate never find a contrastive setting that beats CE only at the
same learning rate by more than about $1\sigma$ (0.0004). When the search is free it switches
the contrastive term off ($w\to0.01$--$0.02$, $\lambda\to0.02$--$0.16$, $\tau\to1$). With $w=1$ and lr $5\cdot10^{-4}$,
14 of 16 contrastive seeds diverge (Section~\ref{sec:stability}).

\begin{table}[htbp]
\centering\small
\caption{GN2-like tagger (Salt), multi-task $\LCE+0.3\,\LJC$, dataset A: best validation CE
(unweighted), mean best epoch, $I_{\mathrm{NCE}}$ at the best epoch ($\log K=8.32$), and test
rejections at 77\% $b$-efficiency (best seed, nominal fractions).}
\label{tab:lambda-gn2}
\begin{tabular}{lcccccc}
\toprule
$\lambda$ & $n$ & best val.\ CE & best epoch & $I_{\mathrm{NCE}}$ [nats] & c-rej & light / $\tau$-rej \\
\midrule
CE only & 8 & $0.60075\pm0.00029$ & 15.6 & --   & 14.45 & 362 / 48.9 \\
0.2     & 8 & $0.60065\pm0.00021$ & 30.6 & 7.90 & 14.57 & 333 / 50.4 \\
0.3     & 6 & $0.60074\pm0.00026$ & 48.7 & 7.67 & 14.56 & 330 / 44.8 \\
0.4     & 5 & $0.60067\pm0.00045$ & 73.2 & 7.26 & 14.50 & 334 / 51.2 \\
0.5     & 5 & $0.60178\pm0.00105$ & 76.2 & 6.50 & 14.49 & 338 / 48.9 \\
0.6     & 5 & $0.60113\pm0.00107$ & 77.6 & 5.89 & 14.58 & 337 / 47.1 \\
0.7     & 4 & $0.60109\pm0.00084$ & 82.5 & 5.47 & 14.45 & 335 / 49.3 \\
\bottomrule
\end{tabular}
\end{table}

\figph{4.5cm}{salt-ririva/slurm/scratch_lambda_scan/output/run_20260928_195621/plots/best_vs_lambda_jets_classification_loss.png}
{GN2-like tagger: best validation CE versus $\lambda$ (same axes as
Fig.~\ref{fig:lambda-fm4tag}; ideally one figure with both models).}{fig:lambda-gn2}

\paragraph{Summary of the 2$\times$2 grid} (best validation CE, unweighted, dataset A):
\begin{center}\small
\begin{tabular}{lcc}
\toprule
 & CE only & CE + JC (best $\lambda$) \\
\midrule
GN2-like (Salt) & $0.60075\pm0.00029$ (8) & $0.60065\pm0.00021$ (8, $\lambda=0.2$) \\
Model A (ours)  & $0.60345\pm0.00049$ (4) & $0.59993\pm0.00078$ (3, $\lambda=0.6$) \\
\bottomrule
\end{tabular}
\end{center}
Without the contrastive term our model is worse than GN2-like; with strong augmentation it
is better. \todo{make the two JC implementations identical, or state the difference: Salt
puts a $128\to256\to128$ projector before the loss, while fm4tag applies $\LJC$ directly to
$z_{\mathrm{jet}}$. The global batch also differs: 4096 (Salt) vs 2048/4096 (fm4tag regimes).}

\subsection{Embedding quality}
At the best epoch:
\begin{itemize}
  \item CE only: uniformity $-1.63$, effective rank $97$;
  \item every contrastive run: uniformity $-3.9$ (the floor), effective rank 121--148.
\end{itemize}
The contrastive term therefore spreads $z_{\mathrm{jet}}$ over the sphere and raises its
dimensionality. Among the contrastive runs, however, the \emph{lower} effective rank goes
with the better CE (Spearman $+0.88$; rank $146\to121$ from $\lambda=0.05$ to $0.8$).
Embedding metrics thus separate ``with and without JC'' but do not rank the JC runs: they
are diagnostics, not targets. During training, the CE-only embedding becomes \emph{less}
uniform (uniformity $-2.43\to-1.20$) while its rank grows ($81\to106$), as the validation CE
overfits (from epoch 7). The $\lambda=0.6$ embedding stays at the uniformity floor
($-3.85$ to $-3.91$) throughout.
\todo{GN2-like runs log the same metrics (Salt \texttt{EmbeddingMetrics}): add them for the
2$\times$2 comparison. Explain the $\lambda=0.5$ rank anomaly (144 in all 5 seeds).}

\figph{4.5cm}{fm4tag/plots/ce_vs_embedding/cl_effective_rank.png}
{Validation CE and effective rank of $z_{\mathrm{jet}}$ per epoch, best CE-only run vs best
contrastive run (companions: \texttt{ce\_only\_effective\_rank.png},
\texttt{*\_uniformity.png}).}{fig:emb-lambda}

\subsection{Flavour-tagging performance}
With the constrained-optimal fractions at 80\% $b$-efficiency:
\begin{itemize}
  \item the best $\lambda$ ($0.55$--$0.65$) reach a $c$-rejection 2.1--2.5\% above the
        GN2-like tagger (10.35) at equal light- and $\tau$-rejection;
  \item CE only is 2.2\% below it;
  \item the statistical error of the test $c$-rejection is $\pm0.27\%$.
\end{itemize}
With the nominal fractions no Model-A tagger beats GN2-like in $c$-rejection. Our models
have much higher light rejection (192--216 vs 171) and lower $\tau$ rejection (23--31 vs 38),
which is why the comparison needs the fraction optimisation. The best validation CE
predicts the optimised $c$-rejection very well (Spearman $-0.98$ over 17 models).
\todo{(a) Re-derive the fractions on validation, not test. (b) Optimise the GN2-like
fractions under the same constraints. (c) Use the multi-seed CE-only cell of the Salt scan
as the reference instead of the single Salt-pt run (weighted CE, 50 epochs, fp16). (d) Give
the seed spread of the $c$-rejection, not only the best seed.}

\figph{5cm}{fm4tag/scripts/plots/run_20261001_110047_lambda_scan_final_wp080_u100_tau100/categorical_no_dropout_min2/ce_only_gn2_best_worst_valce_600dpi/roc_ce_only_gn2_best_second_worst_valce_colours_600dpi.png}
{Light, $c$ and $\tau$ rejection versus $b$-efficiency: CE only, the GN2-like reference,
the best and the worst $\lambda$, each at its constrained-optimal fractions.}{fig:roc-lambda}

% =====================================================================
\section{Results II: contrastive pretraining and training strategies}
\label{sec:res-pretrain}

\note{Purpose: answer Q2 with Model B on dataset B. Every number here is \prov{provisional}
and will be replaced by the WSD redo. Write the text around the qualitative conclusions,
which are the ones expected to survive.}

\subsection{Pretraining}
All three recipes train stably for 30 epochs with QK-norm. For P0 and P1 the validation loss
is still falling at the end (best epoch 28--29), and the training and validation losses
agree (no overfitting on 39.7M jets).
\begin{itemize}
  \item \textbf{P0 and P1}: per-track embeddings reach an effective rank of 114--126 (of 128)
        and uniformity $-3.9$. The jet-feature embedding has rank about 62, with only two
        input features.
  \item \textbf{P1}: the aggregator learns an embedding of rank 130 (of 256).
  \item \textbf{P2} (JC only): the jet-feature embedding collapses to rank 1--2, and in the
        pilot it is barely better than random encoders (Appendix~\ref{app:p2}).
\end{itemize}

\figph{4.5cm}{fm4tag-chunk/plots/pretrain_curves/final/p1_s42_loss_terms.png}
{P1 pretraining: total loss and its terms per epoch, training and validation (companions:
P0, the effective-rank and uniformity versions).}{fig:pretrain}

\subsection{Comparison of training strategies}
Table~\ref{tab:results-ssl}, Fig.~\ref{fig:ssl-summary}:
\begin{itemize}
  \item \textbf{\from{FT}{P1} beats training from scratch} on validation CE
        (\prov{$0.6158\pm0.0002$} vs \prov{$0.6240\pm0.0029$}) and on test $c$-rejection at
        equal light and $\tau$ rejection (\prov{10.8 vs 9.7, +11\%}). With nominal fractions
        it is better on all three rejections.
  \item \textbf{\from{FT}{P0} does not, without dropout} (\prov{0.6312}; with dropout 0.2 it
        becomes the best method, \prov{0.6105} and \prov{+25\%} $c$-rejection, see the note at
        the end of this section): only the per-object embeddings are
        pretrained, and the randomly initialised aggregator is trained at the full learning
        rate. It peaks after 4--5 epochs.
  \item \textbf{Frozen encoders are always worse} than fine-tuning. P1 is better than P0
        for both FT and FR, by 0.015 and 0.014 in CE.
  \item \textbf{The downstream contrastive term does not help.} From scratch, SC+JC = SC
        (\prov{0.6241 vs 0.6240}). Given P1, it hurts: FT \prov{+0.013}, FR \prov{+0.0075}.
\end{itemize}
The production GN2v01 (\prov{$c$-rej 15.6}) is far ahead of every model here, but it was
trained on a much larger, different sample and serves only as a yardstick.
\todo{Fair GN2 reference on dataset B: the Salt GN2-like CE-only run on the 18.6M-jet split
exists (8 seeds, weighted val CE $0.61463\pm0.00062$). Needed: unweighted CE and test
predictions.}

\begin{table}[htbp]
\centering\small
\caption{\prov{Provisional.} Study II, Model B, dataset B: best validation CE (unweighted,
mean $\pm$ std of 4 runs) and test rejections at 77\% $b$-efficiency with the
constrained-optimal fractions (optimised on validation at 80\% against GN2v01) and with the
nominal fractions. $^\dagger$infeasible: the constraints are not met in any run ($f_c=0$).}
\label{tab:results-ssl}
\begin{tabular}{lccccc}
\toprule
Method & best val.\ CE & c-rej (opt.) & light-rej (opt.) & $\tau$-rej (opt.) & c-rej (nominal) \\
\midrule
SC               & \prov{$0.6240\pm0.0029$} & \prov{$9.69\pm0.57$}  & \prov{$434\pm4$}  & \prov{$65.1\pm4.0$} & \prov{$11.99\pm0.45$} \\
SC+JC (clip 1)   & \prov{$0.6241\pm0.0063$} & \prov{$9.79\pm0.92$}  & \prov{$437\pm15$} & \prov{$65.5\pm3.5$} & \prov{$11.97\pm0.48$} \\
\from{FT}{P0}    & \prov{$0.6312\pm0.0007$} & \prov{$8.26\pm0.39$}  & \prov{$452\pm8$}  & \prov{$63.8\pm0.1$} & \prov{$11.05\pm0.28$} \\
\from{FR}{P0}    & \prov{$0.6500\pm0.0008$} & \prov{$4.66\pm0.05^\dagger$} & \prov{$441\pm7$} & \prov{$55.3\pm1.3$} & \prov{$9.71\pm0.09$} \\
\from{FT}{P1}    & \prov{$\mathbf{0.6158\pm0.0002}$} & \prov{$\mathbf{10.78\pm0.10}$} & \prov{$437\pm14$} & \prov{$62.9\pm0.9$} & \prov{$\mathbf{12.62\pm0.06}$} \\
\from{FR}{P1}    & \prov{$0.6361\pm0.0031$} & \prov{$7.42\pm0.41$}  & \prov{$447\pm8$}  & \prov{$65.4\pm1.6$} & \prov{$10.60\pm0.43$} \\
\from{FT}{P1}+JC & \prov{$0.6288\pm0.0002$} & \prov{$8.97\pm0.06$}  & \prov{$442\pm12$} & \prov{$65.0\pm1.3$} & \prov{$11.26\pm0.02$} \\
\from{FR}{P1}+JC & \prov{$0.6436\pm0.0024$} & \prov{$6.41\pm0.65$}  & \prov{$471\pm23$} & \prov{$63.0\pm2.1$} & \prov{$10.05\pm0.22$} \\
\midrule
GN2v01 (yardstick) & -- & 15.58 & 438 & 67.6 & 15.58 \\
\bottomrule
\end{tabular}
\end{table}

\figph{4.5cm}{fm4tag-chunk/plots/arm_comparison/test_metrics/test_rejections_optimal_wp77.png}
{\prov{Provisional.} Test $c$, light and $\tau$ rejections at 77\% for the eight training
methods (4 runs each; companion: best validation CE per method,
\texttt{arm\_comparison/summary/best\_val\_ce\_per\_method.png}).}{fig:ssl-summary}

\note{Current caveats, all addressed by the WSD redo: the batch size differs between
methods (14\,336--94\,208); SC, \from{FT}{P0} and \from{FR}{P0} are early-stopped at
84--95\% of the peak learning rate, while \from{FT}{P1} is still improving when its
40-epoch schedule ends (best epochs 34--39); the P1 CE-only methods reuse learning rates
tuned for P0.}

\subsection{Dependence on the number of labelled jets}
Training on 1, 3, 10, 30 and 100\% of the labels (P0 only, 3 seeds per point;
Table~\ref{tab:labels}, Fig.~\ref{fig:labels}):
\begin{itemize}
  \item P0 pretraining never overtakes training from scratch;
  \item its disadvantage \emph{shrinks} as labels are added (FT: $+0.064$ at 1\% $\to$
        $+0.007$ at 100\%), the opposite of the usual SSL expectation;
  \item every method overfits within a few epochs at $\le10\%$ of the labels, and the
        pretrained ones fastest.
\end{itemize}
Changing the learning rate or adding regularisation moves these numbers by $\le0.01$, much
less than the gaps between methods.

\begin{table}[htbp]
\centering\small
\caption{\prov{Provisional.} Best validation CE versus the fraction of labelled training
jets: mean of 3 downstream seeds (4 at 100\%); below 100\% the pretrained methods start
from one P0 pretraining (seed 42).}
\label{tab:labels}
\begin{tabular}{lrcccc}
\toprule
Labels & jets & SC & \from{FT}{P0} & \from{FR}{P0} & FT$-$SC \\
\midrule
1\%   & 186k  & \prov{$0.7443$} & \prov{$0.8079$} & \prov{$0.8305$} & \prov{$+0.064$} \\
3\%   & 558k  & \prov{$0.7163$} & \prov{$0.7691$} & \prov{$0.7891$} & \prov{$+0.053$} \\
10\%  & 1.86M & \prov{$0.6919$} & \prov{$0.7253$} & \prov{$0.7500$} & \prov{$+0.033$} \\
30\%  & 5.58M & \prov{$0.6640$} & \prov{$0.6759$} & \prov{$0.7064$} & \prov{$+0.012$} \\
100\% & 18.6M & \prov{$0.6240$} & \prov{$0.6312$} & \prov{$0.6500$} & \prov{$+0.007$} \\
\bottomrule
\end{tabular}
\end{table}

\figph{4.5cm}{fm4tag-chunk/plots/label_scan/learning_curve/best_val_ce_vs_labelled_jets.png}
{\prov{Provisional.} Best validation CE versus the number of labelled jets for SC,
\from{FT}{P0} and \from{FR}{P0}.}{fig:labels}

\todo{The label scan is the natural place for a pretraining gain, but it was run only for
P0, the recipe that does not help. Add \from{FT}{P1} (and \from{FR}{P1}) at 1--30\% of the
labels; without it the reader will ask.}

\subsection{Effect of the dataset size}
The same from-scratch recipe trained on dataset A (2.8$\times$ more distinct training jets)
reaches 0.6036, against 0.6240 on dataset B. A batch-size control (2048 instead of 28\,672)
gives 0.6221. The difference in CE between the two studies is therefore due to the amount
of labelled data, not to the model or the batch size.

\note{Dropout study (dropout 0.2, 80 epochs): \from{FT}{P0}+dropout reaches 0.6105 and a
$c$-rejection 25\% above SC, but it is confounded (schedule, batch, untuned lr) and is being
redone. Report it only if it survives the WSD redo, possibly as a short subsection on
regularisation.}

% =====================================================================
\section{Training stability with intersample attention and contrastive objectives}
\label{sec:stability}

\note{Purpose: practical findings that cost weeks and will help others. Keep it to about
one page; it can move to an appendix if space is tight.}

\subsection{Attention-logit growth during contrastive pretraining}
The first 30-epoch pretrainings without QK-norm (lr $10^{-3}$, clip 1.0) failed: 3 of 4
collapsed or spiked. The cause was the row (intersample) attention of the track encoder in
its first blocks:
\begin{itemize}
  \item the scaled logits $q\cdot k/\sqrt d$ grew to $10^5$--$10^7$, so each track attended
        to a single other track of the chunk (normalised entropy $\to0$);
  \item $\|W_{qkv}\|$ grew from 24 to over 300;
  \item bf16 rounding of such logits flips the arg-max and produces gradient spikes.
\end{itemize}
From-scratch supervised models do not show this (logits $\approx20$), so it is specific to
contrastive pretraining. RMS normalisation of queries and keys~\citep{Henry2020QKNorm,Dehghani2023ViT22B}
fixes it: the final P0, P1 and P2 pretrainings (2 seeds each) ran 30 epochs without
spikes. Related observations: \citep{Wortsman2024SmallScaleProxies,Zhai2023EntropyCollapse}.
\todo{the running dropout-0.1 pretrainings show new instabilities with a different
signature (P2 seed 424 collapsed from epoch 0; P0 seed 42 degrading since epoch 9, logits
still about 24). Check them before claiming QK-norm is sufficient.}
\todo{The logit growth has no figure yet; rerun the logit monitor on the old checkpoints.
With a learnable gain, the logits still reach 2.7k--6k in the final P0/P1 without
instability: say ``controlled'', not ``bounded''.}

\subsection{Collapse of multi-task training from scratch}
In Model B, SC+JC without gradient clipping collapsed in 3 of 4 seeds, and with clipping at
5.0 in 2 of 4. Within one epoch the validation CE jumps to the class-prior value (1.29), the
track embeddings collapse to a point (uniformity 0), and the run never recovers. With
clipping at 1.0 none of 4 collapse. This QK-normalised model keeps its attention logits
small (5--9), so the collapse has a different mechanism from the pretraining one. Global
gradient clipping under AdamW does not bound the size of the update.

\subsection{Divergences in the GN2-like tagger}
In Salt, $w=1$ with lr $5\cdot10^{-4}$ diverges in 14 of 16 contrastive seeds. The
divergences are loss spikes at 80--95\% of the peak learning rate. Clipping at 1.0 only
delays them (epoch 20 $\to$ 36). A calibration check (temperature scaling per epoch) shows
that the rise of the validation CE after its minimum is mostly genuine overfitting, not
over-confidence.

% =====================================================================
\section{Discussion}
\label{sec:discussion}

\note{Purpose: interpret, reconcile the two studies, state the limitations honestly.}

\paragraph{The contrastive term as a regulariser.}
In multi-task training, the contrastive term mainly delays overfitting: the CE-only
baselines peak early and overfit, and strong augmentation keeps the validation CE flat for
tens of epochs. This explains why it helps Model A, which overfits strongly (train--validation
gap 0.045 at the early stop), and not the GN2-like tagger. It also explains why a dropout
baseline recovers part of the gain.
\todo{check this interpretation against the regime-pure re-run.}

\paragraph{Why the two studies disagree on JC.}
JC helps Model A on dataset A (52M jets, lr $3\cdot10^{-4}$, long schedules, 2047--4095
negatives) but not Model B on dataset B (18.6M jets, tuned lr $8\cdot10^{-4}$, 40 epochs,
511 negatives per chunk). The paper must either reconcile this with one controlled
experiment or present the two set-ups as distinct. \note{This is the main structural risk
of the paper; see the decision box at the top.}

\paragraph{What makes pretraining work.}
Only pretraining that includes the aggregator (P1) helps, and only when fine-tuned. A
plausible reason is that \from{FT}{P0} starts from a random aggregator trained at the full
learning rate and overfits within 5 epochs, whereas \from{FT}{P1} loads a pretrained
aggregator at 1/10 of the learning rate and keeps improving. \todo{test this directly,
e.g.\ \from{FT}{P0} with the aggregator at lr/10.}

\paragraph{Embedding metrics.}
Uniformity and effective rank detect collapse very well: they mark P2's collapsed
jet-feature embedding, the SC+JC collapses and the pretraining instabilities. They separate runs with
and without a contrastive term. They do \emph{not} rank models by performance (lower rank
is better among contrastive runs).

\paragraph{InfoMin view.}
$\lambda$ controls $I_{\mathrm{NCE}}$ monotonically. A proper InfoMin test (pretraining only,
then a linear probe, performance versus $I_{\mathrm{NCE}}$) was not done. The
frozen-encoder probes of the pilot grid are only a partial proxy.

\paragraph{Limitations.}
\begin{itemize}
  \item Simulation only, one process (\ttbar) and one generator.
  \item The ``unlabelled'' pretraining split has a truth-resampled flavour composition.
  \item The GN2-like reference has no auxiliary tasks.
  \item GN2v01 is not a same-data baseline.
  \item With intersample attention, a jet's prediction depends on its batch-mates: about
        1\% of jets change arg-max when re-chunked. \todo{quantify properly and say how
        inference batches are formed.}
  \item 2--8 seeds per configuration.
  \item Fraction optimisation differs between the studies (to be unified).
\end{itemize}

% =====================================================================
\section{Conclusions and outlook}
\label{sec:conclusions}

\note{Purpose: one paragraph per research question, then the outlook. No new numbers.}

For multi-task training, the augmentation strength decides whether a contrastive term on
the jet embedding helps. With strong augmentation it makes our SAINT-like tagger the best
model on the same jets, slightly ahead of a GN2-like tagger, while it is neutral for GN2
itself. For pretraining, a contrastive objective on the full network, including the
aggregator, followed by fine-tuning, improves the $c$-jet rejection at fixed light and
$\tau$ rejection. Feature-level pretraining and frozen encoders do not help, even with few
labels. Contrastive pretraining with intersample attention requires QK-normalisation.

\emph{Outlook:} robustness to data/MC and MC/MC differences (the original motivation; cf.\
\citep{Sheldon2025MACK,Harris2025Resim}), physics-motivated augmentations (track dropout,
impact-parameter smearing), an InfoMin linear-probe study, pretraining of the GN2-like
tagger in Salt, and scaling of the pretraining sample.

% =====================================================================
\section*{Acknowledgements}
\todo{funding, computing resources (FBK cluster), ATLAS Open Data. Check with the
supervisors and the ATLAS publication policy what an ATLAS member must do to publish a
paper based on ATLAS Open Data (authorship, wording, review, plot labels).}

% =====================================================================
\appendix

\section{Hyper-parameters and compute}
\label{app:hyper}
One table per study with every setting (optimiser, schedule, batch, learning rates,
epochs, early stopping, clipping, dropout, seeds), plus GPU hours.
\begin{itemize}
  \item Pretraining: 34--43\,h per run.
  \item Downstream: 2--18\,h per run.
  \item Model A $\lambda$ cells: up to 88\,h per seed on A40.
\end{itemize}

\section{Pilot studies}
\label{app:pilots}
\begin{itemize}
  \item \textbf{S1, architecture} (scratch CE, 2 seeds): per-token vs concat row attention
        and chunk 128 vs 512 are equal within noise (0.620--0.621). Concat is 1.45$\times$
        faster.
  \item \textbf{S2, pretraining recipe}: grid $\tau\in\{0.1,0.3\}$, $\lambda\in\{0.1,0.4\}$,
        judged by a 15-epoch frozen-encoder probe (random encoders: 0.6563). $\tau=0.1$,
        $\lambda=0.4$ wins for P0 (0.6353), P1 (0.6318) and P2.
  \item \textbf{S3, learning rate per method} (seed 42): grid $\times2$ around
        $3\cdot10^{-4}\sqrt{B/2048}$; for FT, backbone lr = lr (0.643) vs lr/10 (0.630).
\end{itemize}

\section{P2: jet-embedding-only pretraining}
\label{app:p2}
Without per-object terms, the jet-feature embedding receives no direct signal and collapses
(rank 1--2). Unclipped runs collapse entirely (loss $=\ln 511$). The frozen probe beats
random encoders by only 0.005. This negative result motivates P0 and P1, and it contradicts
plan\_v0's ``pretraining on $\LJC$'' as the default recipe.

\section{Fraction optimisation}
\label{app:fractions}
Grid, constraints and feasibility for each model, and the dependence on the working point.

\section{Additional figures}
Per-seed validation curves, rejection versus jet \pT, effective-rank and uniformity curves
for every method, and the SC+JC collapse curves.

% =====================================================================
\bibliographystyle{unsrtnat}
\bibliography{refs}

\end{document}
